\documentclass[UTF8,12pt,a4paper,fontset=fandol]{ctexart}

\usepackage[a4paper,top=25mm,bottom=25mm,left=27mm,right=27mm]{geometry}
\usepackage{amsmath,amssymb,bm}
\usepackage{enumitem}
\usepackage{fancyhdr}
\usepackage{microtype}

\setlength{\parindent}{2em}
\setlength{\parskip}{0.25em}
\setlist[enumerate]{leftmargin=2.2em,itemsep=0.25em,topsep=0.35em}
\allowdisplaybreaks[2]

\pagestyle{fancy}
\setlength{\headheight}{15pt}
\fancyhf{}
\fancyhead[C]{二维对流方程数值测试}
\fancyfoot[C]{\thepage}
\renewcommand{\headrulewidth}{0.4pt}
\renewcommand{\footrulewidth}{0pt}
\makeatletter
\let\ps@plain\ps@fancy
\makeatother

\ctexset{
  section = {
    format = \Large\bfseries,
    number = \arabic{section},
    aftername = \quad
  },
  subsection = {
    format = \large\bfseries,
    number = \thesection.\arabic{subsection},
    aftername = \quad
  },
  subsubsection = {
    format = \normalsize\bfseries,
    number =,
    aftername =
  }
}

\begin{document}

\begin{center}
  {\LARGE\bfseries 第一题：二维对流方程}
\end{center}
\vspace{0.8em}

\section{对流方程}

\subsection{控制方程}

\begin{equation}
\frac{\partial\phi}{\partial t}
+\nabla\cdot(\bm{u}\phi)=0.
\end{equation}

请基于有限体积法写出该方程的半离散形式，并分别采用显式时间离散格式和显式空间离散格式构建数值求解器。

\subsection{数值算例}

\subsubsection{1.2.1\quad 正弦波对流问题}

首先采用下列正弦函数评估线性对流方程的数值精度：

\begin{equation}
\phi_0(x,y)=\sin 2\pi(x+y).
\end{equation}

该正弦波由恒定且均匀的对流速度

\begin{equation}
(u,v)=(1,1)
\end{equation}

进行输运。

计算区域为正方形区域

\begin{equation}
[0,1]^2,
\end{equation}

分别采用三角形单元和四边形单元进行划分，并在 $x$ 和 $y$ 方向施加周期边界条件。

在逐渐加密的网格上进行计算，所有计算中的时间步长均根据

\begin{equation}
\mathrm{CFL}=0.2
\end{equation}

进行调整。

请给出 $t=1.0$ 时的 $L_1$ 误差和数值收敛阶。

\subsubsection{1.2.2\quad 复杂轮廓的刚体旋转问题}
\thispagestyle{fancy}

为了评估数值格式分辨不连续结构的能力，计算由开槽圆盘、圆锥和光滑凸峰组成的复杂轮廓的刚体旋转问题。

旋转速度场为

\begin{equation}
\bm{u}=(0.5-y,\;x-0.5).
\end{equation}

计算区域为正方形区域

\begin{equation}
[0,1]^2.
\end{equation}

初始标量场由开槽圆盘、圆锥和光滑凸峰组成：

\begin{equation}
\phi_0(x,y)
=\phi_D(x,y)+\phi_C(x,y)+\phi_H(x,y).
\end{equation}

三个轮廓的半径均为

\begin{equation}
r_0=0.15.
\end{equation}

\paragraph{（1）开槽圆盘}

开槽圆盘的圆心为

\begin{equation}
(x_D,y_D)=(0.5,0.75).
\end{equation}

定义

\begin{equation}
r_D(x,y)
=\sqrt{(x-0.5)^2+(y-0.75)^2}.
\end{equation}

开槽圆盘为

\begin{equation}
\phi_D(x,y)
=
\begin{cases}
1,
&
r_D(x,y)\leq 0.15
\quad\text{且}\quad
\left(
|x-0.5|\geq 0.025
\ \text{或}\
y\geq 0.85
\right),\\[2mm]
0,
&\text{其他}.
\end{cases}
\end{equation}

\newpage
\paragraph{（2）圆锥}

圆锥的中心为

\begin{equation}
(x_C,y_C)=(0.5,0.25).
\end{equation}

定义

\begin{equation}
r_C(x,y)
=\sqrt{(x-0.5)^2+(y-0.25)^2}.
\end{equation}

圆锥为

\begin{equation}
\phi_C(x,y)
=
\begin{cases}
1-\dfrac{r_C(x,y)}{0.15},
&r_C(x,y)\leq0.15,\\[3mm]
0,
&r_C(x,y)>0.15.
\end{cases}
\end{equation}

\paragraph{（3）光滑凸峰}

光滑凸峰的中心为

\begin{equation}
(x_H,y_H)=(0.25,0.5).
\end{equation}

定义

\begin{equation}
r_H(x,y)
=\sqrt{(x-0.25)^2+(y-0.5)^2}.
\end{equation}

光滑凸峰为

\begin{equation}
\phi_H(x,y)
=
\begin{cases}
\dfrac14
\left[
1+
\cos\left(
\pi\dfrac{r_H(x,y)}{0.15}
\right)
\right],
&r_H(x,y)\leq0.15,\\[4mm]
0,
&r_H(x,y)>0.15.
\end{cases}
\end{equation}

计算区域分别采用三角形单元和四边形单元进行划分，每条边布置 $100$ 个顶点。

请绘制完成一周旋转后，即

\begin{equation}
t=2\pi
\end{equation}

时的标量场等值线图。

\end{document}
